\subsection{A singleton stalk opposite a conical block}

\begin{proposition}[Proved: an extreme fixed arrow opposite a conical block]
\label{CP2:A-singleton-stalk-conical-block}
Let $C$ be a finite category with acyclic support and
let $p$ be an extreme object. Suppose the full subcategory
$E=C\setminus\{p\}$ has an initial or terminal object.
If at least one nonempty strict hom-set incident with
$p$ is a singleton, every rational degree-two product
on $BC$ vanishes. There is no restriction on the order
of the corner monoid at $p$ or on the other stalk sizes.
\end{proposition}
\begin{proof}
Write $M=\operatorname{End}_C(p)$. Since $p$ is extreme,
$C$ is a one-way collage of $M$ and $E$. The space
$BE$ is contractible because $E$ has an initial or
terminal object. In particular $H^2(BE;\mathbb Q)=0$.
Let $a$ be the unique arrow in the specified singleton
stalk. If $p$ is a source, every right $M$ actor fixes
$a$; if $p$ is a sink, every left $M$ actor fixes $a$.
The fixed arrow gives a section $BM\to T$ of the
bridge projection in the corresponding two-block
homotopy-pushout model. Consequently the pullback
$H^2(BM;\mathbb Q)\to H^2(T;\mathbb Q)$ is injective.

The Mayer--Vietoris sequence now makes the restriction
of each second cohomology class of $BC$ to the two
vertex blocks zero. Equivalently, apply
Corollary~\ref{CP2:A-fixed-arrow} or its opposite.
Each class lifts from relative cohomology of those
blocks, and the relative quotient is a suspension by
Lemma~\ref{CP2:A-two-block}. Products of these classes
are therefore zero, as claimed.
\end{proof}

\begin{corollary}[Proved: the frozen four-object profile $2285$]
\label{CP2:A-four-eighteen-profile-2285}
The eighteen-arrow profile with hom-set matrix
\[
\begin{pmatrix}
1&0&1&1\\
0&7&2&1\\
0&0&1&3\\
0&0&0&1
\end{pmatrix}
\]
cannot have classifying-space type $\CP^2$.
\end{corollary}
\begin{proof}
Object $1$ is a source and its stalk to object $3$ is
a singleton. Its deletion block on $\{0,2,3\}$ has
initial object $0$: there is exactly one arrow from
$0$ to each of those objects. Apply the proposition.
\end{proof}
